\documentclass[10pt]{article}

\usepackage[utf8]{inputenc}

\usepackage[T1]{fontenc}

\usepackage{amsmath}

\usepackage{amsfonts}

\usepackage{amssymb}

\usepackage[version=4]{mhchem}

\usepackage{stmaryrd}

\usepackage{graphicx}

\usepackage[export]{adjustbox}

\graphicspath{ {./images/} }



\begin{document}

Часто в этом классе методов могут быть найдены такие, что 1) весь переход от сеточного слоя в момент времени $t_{n}$ к новому сеточному слою $t=t_{n+1}$ является достаточно простым и может быть осуществлен при затрате $O(N)$ арифметич. действий, где $N$ - число узлов пространственной сетки; 2) гарантируется абсолютная устойчивость метода, 3) гарантируется наличие приемлемой точности метода (наличие аппроксимации в том или ином смысле). Р. м. довольно широко применяются при практич. решении многомерных задач математич. физики, связанных, напр., с линейными и нелинейными системами параболического, гиперболического или смешанного типа (см. [1] - [8]).



Обычно для задачи с $p$ пространственными переменными переход от $t_{n}$ к $t_{n+1}$ в Р. м. производится с исшользованием $p$ вспомогательных (дробных) шагов:



\[

A_{\mathrm{S}}^{(n)} u^{n+s / p}=B_{\mathrm{S}}^{(n)}\left(u^{n}, u^{n+1 / p}, \ldots, u^{n+(s-1) / p}\right),

\]



где



\[

u^{n} \equiv u\left(t_{n}\right), u^{n+k / p} \equiv u\left(t_{n+1}\right), A_{\mathrm{s}}^{(n)}

\]



\begin{itemize}

  \item матрица, соответствующая разностной аппроксимации нек-рого диффференциального оператора, содержащего производные только по $x_{s}$ (одномерного дифференциального оператора), а правые части (1) легко вычислимы. При соответствующей нумерации неизвестных, связанной с выбором направления $x_{s}$, матрицы $A_{s}^{(n)}$ становятся обычно диагональными и решение систем (1) при каждом $s$ сводится к многократному решению одномерных разностных систем по направлению $x_{s}$. Поэтому часто Р. м. наз. также или переменных направлений методом или дробных шагов методом.

\end{itemize}



Одним из типичных примеров в случае уравнения



\[

\frac{\partial u}{\partial t}=\frac{\partial^{2} u}{\partial x_{1}^{2}}+\frac{\partial^{2} u}{\partial x_{2}^{2}}+f, \quad\left(x_{1}, x_{2}\right) \in \Omega,

\]



с начальным условием $\left.u\right|_{t=0}=\varphi\left(x_{1}, x_{2}\right)$ и краевым условием $\left.u\right|_{\Gamma}=\psi\left(x_{1}, x_{2}, t\right)$, где $\Omega \equiv(0,1) \times(0,1), \Gamma-$ граница $\Omega$, может служить следующий метод, построенный на квадратной сетке с шагом $h$ :



\[

\left.\begin{array}{cc}

\frac{u_{i}^{n+1 / 2}-u_{i}^{n}}{\tau / 2}+\left[\Lambda_{1}\left(\sigma u^{n+1 / 2}+(1-\sigma) u^{n}\right)\right]_{i}=f_{i}^{n+1 / 2}, \\

\overline{x_{i}} \in \Omega_{n} ; \\

u_{i}^{n+1 / 2}=\psi_{i}^{n+1 / 2} \text { при } \bar{x}_{i} \in \Gamma, u_{i}^{n+1}=\psi_{i}^{n+1} \text { при } \\

\frac{u_{i}^{n+1}-u_{i}^{n+1 / 2}}{\tau / 2}+\left[\Lambda_{2}\left(\sigma u^{n+1}+(1-\sigma) u^{n+1 / 2}\right)\right]_{i}= \\

=f_{i}^{n+1 / 2}, & \overline{x_{i}} \in \Omega_{n},

\end{array}\right\}

\]



где $i \equiv\left(i_{1}, i_{2}\right), \bar{x}_{i} \equiv\left(i_{1} h, i_{2} h\right), u_{i}^{n} \equiv u\left(n \tau, \bar{x}_{i}\right), \psi_{i}^{n+1 / 2} \equiv$ $\equiv \psi\left((n+1 / 2) \tau, \bar{x}_{i}\right),-\Lambda_{1}$ и $\Lambda_{2}$ - простейшие разностные аппроксимации для $\frac{\partial^{2}}{\partial x_{1}^{2}}$ и $\frac{\partial^{2}}{\partial x_{2}^{2}}, \Omega_{n}-$ совокупность внутренних узлов $\bar{x}_{i}, \sigma \geqslant 1 / 2$.



Имеются два альтернативных подхода к теории Р. м. В одном из них промежуточные шаги ни в чем существенном не отличаются от целых шагов, и сами разностные уравнения на дробных шагах и граничные условия для них, подобно методу (2), устроены одинаково, и можно ожидать, что $u^{n}$ и $u^{n+1 / 2}$ будут служить аппроксимациями для решения исходной задачи в моменты времени $t_{n}$ и $t_{n+1 / 2} \equiv t_{n}+\tau / 2$. Этот подход основан на использовании понятия составной схемы и суммарной аппроксимации (см. [2]). Схемы такого типа часто наз. ло к а ль но о дноме риым и с х е ма и или а д д и т и в ны м и с х е м а м и; их можно также трактовать как обычные разностные схемы для некрого уравнения с сильно осциллирующими по времени коәффициентами, решение к-рого должно быть близко к решению исходной задачи (см. [1]-[4]). Достоинства этого подхода в его простоте и общности, напр. обобпееия метода (2) возможны и для случая криволинейных областей $\Omega$ и более общих задач. Точность же получаемых на этом пути методов обычно не очень высока. Известны и иногда успешно применяются варианты Р. м., в к-рых расщепление производится не по пространственным переменным, a по физич. процессам (cм. [5]).



Второй подход в плане анализа устойчивости и сходимости исключает какие-либо дробные шаги из рассмотрения. Сама разностная схема и аппроксимация трактуются традиционным образом. Необычность разностной схемы проявляется лишь в том, что на верхнем слое схемы появляется необычный разностный оператор. Напр., вместо метода (2) рассматривается метод



\[

\begin{gathered}

\left(A\left(\frac{u^{n+1}-u^{n}}{\tau}\right)\right)_{i}+\left(\Lambda_{1} u^{n}\right)_{i}+\left(\Lambda_{2} u^{n}\right)_{i}=f_{i}^{n+1 / 2}, \\

\bar{x}_{i} \in \Omega_{n} \\

u_{i}^{n}=\psi_{i}^{n} \quad \text { при } \bar{x}_{i} \in \Gamma,

\end{gathered}

\]



где $A \equiv\left(E+\sigma \tau \Lambda_{1}\right)\left(E+\sigma \tau \Lambda_{2}\right), E$ - тождественный оператор. Такие операторы $A$ обычно наз. р а с щ е п л яющим и с я или фа т о р и зов а н ны м и о п е р а т о р а м и. Дробные шаги связываются лишь с методом решения возникающих систем и для одной и той же схемы (3) могут быть введены различными способами, граничные условия для них должны выбираться в зависимости от этого. Сами схемы типа (3) можно трактовать как обычные схемы с весом для $\varepsilon$-уравнения, напр., вида



\[

\frac{\partial u}{\partial t}-\frac{\partial^{2} u}{\partial x_{1}^{2}}-\frac{\partial^{2} u}{\partial x_{2}^{2}}+\varepsilon \frac{\partial^{5} u}{\partial x_{1}^{2} \partial x_{2}^{2} \partial t}=f, \quad \varepsilon>0,

\]



решения к-рого отличаются от решения исходной задачи на $O(\varepsilon)$ (см. [4]). В случаз области $\Omega$, составленной из прямоугольников, матрицы возникающих систем в методах типа (3) уже не представимы в виде произведения "одномерных" матриц. Все же решения подобных систем могут быть найдены при затрате $O(N)$ арифметич. действий (см. [4]), операторы подобного типа наз. р а с ши р енно р а с щ п л я ю пи мися опер а т о р а м и. При исследовании устойчивости и сходимости схем с расщепляющимися и расширенно расщепляющимися операторами большую роль играет метод энергетич. неравенств (см. [2], [4], [6] - [8]).



Лит.: [1] м а р ч у к Г. И., Методы вычислительной мате матики, 2 изд., М., 1980 ; [2] С а м а р с к и й А. А., Теория разностных схем, М., 1977; [3] Я н е н к о Н. Н., Метод дробных шагов решения многомерных задач математической физики, Новосиб., 1967; [4] Д ь я к о н о в Е. Г., Разностные методы решения краевых задач, в. 1-2, М., 1971-72; [5] К о в е н я В. М., Я н е н к о Н. Н., Метод расщепления в задачах газовой динамики, Новосиб., 1981, [6] Д. р ы я М., в кн.. Banach center pul



\includegraphics[max width=\textwidth, center]{2023_07_08_b12a30ddd1d69094bec0g-1(1)}

[8] H a y e s L. J., "SIAM J. Numer. Analysis", 1981 , v. 18, № 4



\begin{center}

\includegraphics[max width=\textwidth]{2023_07_08_b12a30ddd1d69094bec0g-1}

\end{center}



РАСІІЕПЛЯЕМАЯ ГРУППА - группа $G$, порожденная своими подгруппами $H$ и $K$ такими, что $H$ инвариантна в $G$ и пересечение $H \cap K=E$ (так что факторгруппа $G / H$ изоморфна $K$ ). В этом случае $G$ наз. также р а сщ е пл я е мым р а сши р е ни е м группы $H$ при помощи группы $K$, или по л у пр ямы м п р о и з в е д е н и е м $H$ на $K$. Если подгруппы $H$ и $K$ поэлементно перестановочны, то их полупрямое произведение совпадает с прямым произведением $H \times K$. ІІолупрямое произведение $G$ группы $H$ на группу





\end{document}